\chapter{Time and Temporal Facts}
\label{ch:temporal}

\Morendo\ has three levels of support for time: functions that compare
and shift instants, facts that carry effective and expiration times, and
a conditional element that only matches fresh facts. This chapter covers
them in that order; rule-level dates were covered in
Section~\ref{sec:rule-dates}.

Temporal logic in general reasons with words such as \emph{always},
\emph{never}, \emph{sometimes} and \emph{eventually}, and an engine for
it would need an agreed vocabulary of such terms, from \emph{weekly} and
\emph{quarterly} to \emph{immediately}. \Morendo\ does not attempt that.
It supplies the practical pieces below and defines no vocabulary of its
own; an application that needs a term such as ``business day'' defines
it with functions.

\section{Time values and functions}

A time is a \fn{java.time.Instant}. \fn{(now)} returns one, and a slot
declared \fn{(type DATE)} is meant to hold one; the time functions also
accept a \fn{java.util.Date} or \fn{Calendar} coming from a Java bean and
a plain number of epoch milliseconds, so log data with numeric
timestamps works without conversion. Times print in ISO-8601 form.

\begin{lstlisting}[style=shell]
Morendo> (bind ?t1 (now))
true
Morendo> (bind ?t2 (add-hours 2 ?t1))
true
Morendo> (before ?t1 ?t2)
true
Morendo> (within-minutes 30 ?t1 ?t2)
false
Morendo> (eq-day ?t1 ?t2)
true
\end{lstlisting}

The full set is in Appendix~\ref{app:functions}: \fn{add-seconds},
\fn{add-minutes}, \fn{add-hours}; \fn{before}, \fn{after},
\fn{between}; \fn{eq-second} through \fn{eq-year}, which compare after
truncation; and \fn{within-seconds} through \fn{within-days}, which
test the distance between two times. In a rule they appear in predicate
constraints, which is how the MOLAP sample joins two logs that happened
within a few minutes of each other:

\begin{lstlisting}
(DatabaseQueryLog
  (endTime ?dbEndTime&:(within-minutes 4 ?dbEndTime ?endTime))
  (timestamp ?dbTimestamp&:(within-minutes 2 ?timestamp ?dbTimestamp)))
\end{lstlisting}

\file{samples/time/} has three one-line examples.

\section{Temporal facts}
\label{sec:temporal-facts}

A temporal fact is a fact with a validity interval. It is asserted with
\fn{assert-temporal}, which accepts, besides the template's own slots,
the pseudo-slots \fn{effective-time} and \fn{expiration-time} (epoch
milliseconds), \fn{source}, \fn{service-type} and \fn{validity}:

\begin{lstlisting}[style=shell]
Morendo> (deftemplate evt (slot a))
true
Morendo> (assert-temporal (evt (effective-time 0) (expiration-time 4102444800000) (a "valid")))
2
Morendo> (facts)
f-1 (_initialFact)
f-2 (evt (a "valid") (effective-time 0)(expiration-time 4102444800000)(service-type )(source )(validity 0))
for a total of 2
\end{lstlisting}

From Java, \fn{assertTemporalObject} and the temporal \fn{assertFact}
take \fn{Instant}s (Chapter~\ref{ch:embedding}). Temporal facts match
patterns like any other fact. A rule declared
\fn{(temporal-activation TRUE)} checks the facts of each match when it
reaches the rule's terminal node: if one of them is a temporal fact whose
expiration time has passed, the activation is not added and the expired
facts are retracted as soon as the current assertion has finished
propagating.

\begin{lstlisting}[style=shell]
Morendo> (reset)
Morendo> (defrule ta (declare (temporal-activation TRUE)) (evt (a ?x))
  => (printout t "ta fired " ?x crlf))
true
Morendo> (assert-temporal (evt (effective-time 0) (expiration-time 1) (a "expired")))
2
Morendo> (assert-temporal (evt (effective-time 0) (expiration-time 4102444800000) (a "valid")))
3
Morendo> (facts)
f-1 (_initialFact)
f-3 (evt (a "valid") (effective-time 0)(expiration-time 4102444800000)(service-type )(source )(validity 0))
for a total of 2
Morendo> (fire)
ta fired valid
1
\end{lstlisting}

Rules without the property see expired facts like any other, so a rule
set that mixes both kinds decides per rule whether expiry matters.

\section{The temporal conditional element}

Some conditions hold only for a while. ``If the customer is in Boston,
send a coupon'' is true while a driver from Rhode Island to Maine passes
through Boston, and false before and after; a rule written with ordinary
patterns would go on matching the old position until something
retracted it. The \fn{temporal} element is a pattern with a freshness
window for such cases:

\begin{lstlisting}
(temporal (declare ?tv (relative-time <seconds>)
                   [(interval-time <seconds> [(<function> <number>)])])
          (<template> (<slot> <constraint>)...))
\end{lstlisting}

It is compiled into a temporal join, and the window applies to the facts
of the temporal pattern itself. When a combination of facts from the
preceding elements reaches the join, it is combined only with facts of
the temporal pattern asserted less than \fn{relative-time} seconds ago;
older ones no longer match, and the join drops them from its memory as it
meets them. A fact of the temporal pattern that arrives later is combined
at once with the waiting facts of the preceding elements, however old
they are, unless the element before it is temporal too, in which case
that element's window applies to them. Either way the facts stay in
working memory until something retracts them.

\begin{lstlisting}
(defrule match_order
  (transaction (txType "sell") (shares ?shares) (sym ?sym))
  (temporal (declare ?tx (relative-time 300))
    (transaction (exchange "NYSE") (txType "buy") (shares ?shares) (sym ?sym)))
=>
  (printout t "sell of " ?shares " " ?sym " matches a recent buy" crlf))
\end{lstlisting}

Here a new sell order is paired with buy orders for the same stock and
number of shares asserted in the last five minutes; a buy order older
than that is not considered.

\begin{lstlisting}
(defrule pair
  (evt (a ?x) (k ?k))
  (temporal (declare ?tv (relative-time 10)) (evt (a ?y) (k ?k)))
=>
  (printout t "pair " ?x " " ?y crlf))
\end{lstlisting}

The window is measured from the time a fact was asserted, not from a
slot value. The join is keyed by the variables the temporal pattern
shares with the earlier patterns (\fn{?k} above); with no shared
variable every combination within the window qualifies, including a
fact with itself. Predicate constraints inside a temporal pattern take
part in that key rather than being evaluated, so put conditions that
relate the two facts, such as \fn{(neq ?y ?x)}, in a \fn{test} element
after the temporal one. There is no sample using the element in the
repository.

With \fn{interval-time} the join does not pass combinations on as they
form. It collects them and releases them in a batch once the interval,
in seconds, has elapsed; there is no timer, so a batch goes out when the
next fact reaches the join after the interval is over. The call inside
the declaration names a Java function that receives the number and the
list of collected matches, and returns the matches to pass on; without
it every collected match is passed on. A match whose facts are retracted
while it waits is dropped.

\section{Temporal distance of templates}

\fn{(set-temporal-distance template seconds)} records for a template how
long its facts remain relevant, and \fn{(calculate-temporal-distance file
out)} derives the distance of every template used by the rules of the
current module from their \fn{temporal} elements and writes a report.
These are bookkeeping for temporal rule sets and have no effect on
matching.

\section{Practical advice}

For event streams the combination that works today is: a
\fn{no-agenda} rule that reacts to each event as it arrives
(Section~\ref{sec:no-agenda}), time functions in its predicates to
relate events, and a housekeeping rule or a Java timer that retracts old
facts by their timestamp slot. Keep timestamps as instants or epoch
milliseconds in ordinary slots so that they can be compared with
\fn{before}, \fn{within-seconds} and friends.
